% Titolo: Mosaico e pianificazione sprint
\documentclass[a4paper, 11pt]{article}

% --- PACCHETTI ---
\usepackage[utf8]{inputenc}
\usepackage[italian]{babel}

\usepackage[margin=2.5cm]{geometry}
\usepackage{graphicx}
\usepackage{tabularx}
\usepackage{hyperref}
\usepackage{fancyhdr}
\usepackage{array}
\usepackage{titlesec}
\usepackage{xcolor}

% --- CONFIGURAZIONE HEADER E FOOTER ---
\pagestyle{fancy}
\fancyhf{}
\fancyhead[L]{Bit Crusher}
\fancyhead[R]{Verbale Interno N. 9}
\fancyfoot[C]{\thepage}
\renewcommand{\headrulewidth}{0.4pt}

% Configurazione link indice
\hypersetup{
    colorlinks=true,
    linkcolor=black,
    urlcolor=blue
}

\begin{document}

% ==========================================
% FRONTESPIZIO
% ==========================================
\begin{titlepage}
    \centering
    
    % --- LOGHI ---
    % Sostituisci i \framebox con i comandi \includegraphics corretti
    \begin{minipage}{0.45\textwidth}
        \centering
        % Esempio: \includegraphics[width=0.7\textwidth]{logo_unipd.png}
        \includegraphics[width=0.7\textwidth]{unipd_logo.png}
    \end{minipage}
    \hfill
    \begin{minipage}{0.45\textwidth}
        \centering
        % Esempio: \includegraphics[width=0.7\textwidth]{WhatsApp Image 2026-10-15 at 17.29.12.jpeg}
        \includegraphics[width=0.7\textwidth]{Logo-BitCrusher-Nero.png}
    \end{minipage}

    \vspace{1.5cm}
    {\Large \textbf{UNIVERSITÀ DEGLI STUDI DI PADOVA}} \\
    \vspace{0.3cm}
    {\large Corso di Laurea in Informatica} \\
    \vspace{0.3cm}
    {\large Insegnamento di Ingegneria del Software} \\
    \vspace{0.3cm}
    {\large Anno Accademico 2026/2027} \\

    \vspace{2cm}
    {\Huge \textbf{Verbale di Riunione}} \\
    \vspace{0.5cm}
    {\Large Riunione Interna \#9}

    \vspace{2cm}
    \begin{flushleft}
        \large
        \textbf{Gruppo:} Bit Crusher \\
        \textbf{Email:} bitcrusherteam@gmail.com \\
        \textbf{Versione:} v2.0.0 \\
        \textbf{Data della riunione:} 2026-11-16 \\
        \textbf{Data ultima modifica:} 26 novembre 2026
    \end{flushleft}
    \vfill
\end{titlepage}

\newpage

% ==========================================
% REGISTRO DELLE MODIFICHE
% ==========================================
\section*{Registro delle Modifiche}
\vspace{0.5cm}
\begin{center}
\renewcommand{\arraystretch}{1.5}
\begin{tabularx}{\textwidth}{|l|l|l|l|X|}
\hline
\textbf{Versione} & \textbf{Data} & \textbf{Autore} & \textbf{Verificatore} & \textbf{Descrizione} \\ \hline
2.0.0 & 2026-11-26 & Erik Giuliano & - & Approvazione per il rilascio. \\ \hline
1.0.1 & 2026-11-26 & Erik Giuliano & - & [che cosa cambia: da scrivere] \\ \hline
1.0.0 & 2026-10-21 & Erik Giuliano & - & Approvazione per il
rilascio. \\ \hline
0.1.0 & 2026-10-21 & Riccardo Pozzato & Antonio Tufano & Punto della situazione sui vecchi To-Do e creazione nuovi To-Do. \\ \hline

\end{tabularx}
\end{center}

\newpage

% ==========================================
% INDICE
% ==========================================
\tableofcontents
\newpage
% ==========================================
% 1. INFORMAZIONI GENERALI
% ==========================================
\section{Informazioni Generali}
\begin{itemize}
    \item \textbf{Luogo/Piattaforma:} Riunione Discord
    \item \textbf{Data:} 2026-11-16
    \item \textbf{Orario di inizio:} 17:30
    \item \textbf{Orario di fine:} 19:15
    \item \textbf{Responsabile:} Erik Giuliano
    \item \textbf{Scriba:} Riccardo Pozzato
\end{itemize}

\vspace{0.5cm}
\textbf{Partecipanti presenti:} \\
\vspace{0.2cm}
\renewcommand{\arraystretch}{1.2}
\begin{tabularx}{0.8\textwidth}{|X|c|}
\hline
\textbf{Nome e Cognome} & \textbf{Matricola} \\ \hline
Pietro Leonardo Acampora & 2146762 \\ \hline
Marco Bollettin & 2147955 \\ \hline
Riccardo Pozzato & 2147988 \\ \hline
Alex Singh & 2148621 \\ \hline
Antonio Tufano & 2137972 \\ \hline
Erik Giuliano & 2138012 \\ \hline
\end{tabularx}

\vspace{0.5cm}
\textbf{Partecipanti assenti:} \\
\vspace{0.2cm}
\renewcommand{\arraystretch}{1.2}
\begin{tabularx}{0.8\textwidth}{|X|c|}
\hline
\textbf{Nome e Cognome} & \textbf{Matricola} \\ \hline
 Alexandru Mitu & 2101083 \\ \hline
\end{tabularx}

% ==========================================
% 2. ORDINE DEL GIORNO
% ==========================================
\vspace{0.5cm}
\section{Ordine del Giorno}
Gli argomenti previsti per la riunione odierna sono:
\begin{enumerate}
    \item Verifica dello stato di svolgimento dei vecchi To-Do.
    \item Creazione ed assegnazione nuovi To-Do.
\end{enumerate}

% ==========================================
% 3. RESOCONTO DELLA RIUNIONE
% ==========================================
\newpage
\section{Resoconto della Riunione}

\subsection{Verifica dello stato di svolgimento dei vecchi To-Do}
\begin{itemize}
    \item OP\#42 (redazione dei contenuti preliminari del glossario strutturato): chiusa.
    \item OP\#43 (definizione del calendario degli incontri con Var Group): chiusa.
    \item OP\#44 (stesura delle Norme di Progetto): in corso.
    \item OP\#45 (automatizzazione della documentazione): in corso; mancano il rendering del PDF, la configurazione del bot di GitHub per l'invio diretto dei commit e il runner per la compilazione del codice LaTeX.
    \item OP\#46 (miglioramento della tabella di tracciamento): aperta; l'attività è stata completata, ma le modifiche non sono ancora state caricate sul repository in attesa della risoluzione di un'attività bloccante.
    \item OP\#47 (valutazione del calendario condiviso): aperta.
    \item OP\#48 (sincronizzazione delle pull request di GitHub con la vista a mosaico): chiusa.
    \item OP\#49 (definizione di una struttura compatibile con il sorgente LaTeX del glossario): chiusa.
    \item OP\#52 (approvazione del quarto diario di bordo): chiusa.
\end{itemize}


% ==========================================
% 4. AZIONI DA SVOLGERE (TO-DO)
% ==========================================
\vspace{0.5cm}
\section{Azioni da Svolgere (To-Do)}
\begin{tabularx}{\textwidth}{|c|X|l|l|}
\hline
\textbf{Codice Rif.} & \textbf{Descrizione Task} & \textbf{Assegnatario} & \textbf{Scadenza} \\ \hline
OP\#53 & Redazione del Diario di Bordo \#5 & Marco Bolletin
 & 2026-11-20 \\ \hline
OP\#54 & Approvazione quinto Diario di Bordo & Erik Giuliano
 & 2026-11-23 \\ \hline
OP\#55 & Redazione del Verbale Interno \#9 & Riccardo Pozzato
 & 2026-11-19 \\ \hline
OP\#56 & Verifica del Verbale Interno \#9 & Antonio Tufano
 & 2026-11-20 \\ \hline
OP\#57 & Perimetrare gli artefatti sulla repo da OpenProject & Antonio Tufano
 & 2026-11-17 \\ \hline
OP\#58 & Rendere la vista web più fruibile & Pietro L. Acampora
 & 2026-11-20 \\ \hline
\end{tabularx}

% ==========================================
% 5. PROSSIMA RIUNIONE
% ==========================================
\vspace{0.5cm}
\section{Prossima Riunione}
\begin{itemize}
    \item \textbf{Data:}  2026-11-19
    \item \textbf{Ora:} 17:30
    \item \textbf{OdG preliminare:} Definizione delle convenzioni per i verbali (formato delle date, codici univoci e regole di nomenclatura).
\end{itemize}

\end{document}